
\documentclass{article}
\usepackage{colortbl}
\usepackage{makecell}
\usepackage{multirow}
\usepackage{supertabular}

\begin{document}

\newcounter{utterance}

\centering \large Interaction Transcript for game `cladder', experiment `full\_v1.5\_default', episode 7073 with qwen3.5{-}4b.
\vspace{24pt}

{ \footnotesize  \setcounter{utterance}{1}
\setlength{\tabcolsep}{0pt}
\begin{supertabular}{c@{$\;$}|p{.15\linewidth}@{}p{.15\linewidth}p{.15\linewidth}p{.15\linewidth}p{.15\linewidth}p{.15\linewidth}}
    \# & \multicolumn{2}{c}{Player} && \multicolumn{2}{c}{Game Master} \\
    \hline

    \theutterance \stepcounter{utterance}  
    & & & \multicolumn{4}{p{0.6\linewidth}}{
        \cellcolor[rgb]{0.9,0.9,0.9}{
            \makecell[{{p{\linewidth}}}]{
                \texttt{\tiny{[P1$\langle$GM]}}
                \texttt{You are an expert in causal inference. The following question is not a typical commonsense query, but rather a meticulously designed question created by a professor specializing in causal inference, intended to assess the students' mastery of the course content.} \\
\texttt{Imagine a self{-}contained, hypothetical world with only the following conditions, and without any unmentioned factors or causal relationships: Jyka has a direct effect on yupt. Yupt has a direct effect on kwox. The overall probability of jyka is 25\%. The probability of not jyka and kwox is 71\%. The probability of jyka and kwox is 24\%. Is the chance of kwox smaller when observing jyka?} \\
\texttt{Start your answer with 'Yes' or 'No', followed by additional reasoning or evidence to support your explanation.} \\
            }
        }
    }
     \\ \\

    \theutterance \stepcounter{utterance}  
    & \multicolumn{4}{p{0.6\linewidth}}{
        \cellcolor[rgb]{0.561,0.737,0.561}{
            \makecell[{{p{\linewidth}}}]{
                \texttt{\tiny{[P1$\rangle$GM]}}
                \texttt{Yes} \\
\\ 
\texttt{**Reasoning:**} \\
\\ 
\texttt{To determine whether the chance of *kwox* is smaller when observing *jyka*, we must compare the marginal probability of *kwox* (\$P(\textbackslash{}text\{kwox\})\$) with the conditional probability of *kwox* given *jyka* (\$P(\textbackslash{}text\{kwox\} | \textbackslash{}text\{jyka\})\$).} \\
\\ 
\texttt{1.  **Calculate the marginal probability of *kwox* (\$P(\textbackslash{}text\{kwox\})\$):**} \\
\texttt{    We are given the probability of the joint event where *jyka* is false and *kwox* is false: \$P(\textbackslash{}neg \textbackslash{}text\{jyka\} \textbackslash{}cap \textbackslash{}neg \textbackslash{}text\{kwox\}) = 0.71\$.} \\
\texttt{    We are also given the probability of the joint event where *jyka* is true and *kwox* is true: \$P(\textbackslash{}text\{jyka\} \textbackslash{}cap \textbackslash{}text\{kwox\}) = 0.24\$.} \\
\\ 
\texttt{    The sum of these two disjoint events represents the probability that *kwox* is false:} \\
\texttt{    \$P(\textbackslash{}neg \textbackslash{}text\{kwox\}) = P(\textbackslash{}neg \textbackslash{}text\{jyka\} \textbackslash{}cap \textbackslash{}neg \textbackslash{}text\{kwox\}) + P(\textbackslash{}text\{jyka\} \textbackslash{}cap \textbackslash{}neg \textbackslash{}text\{kwox\})\$} \\
\\ 
\texttt{    However, we can more directly find \$P(\textbackslash{}text\{kwox\})\$ by subtracting the probability of *kwox* being false from 1. To do this, we first need \$P(\textbackslash{}neg \textbackslash{}text\{kwox\})\$.} \\
\texttt{    We know \$P(\textbackslash{}neg \textbackslash{}text\{jyka\}) = 1 {-} P(\textbackslash{}text\{jyka\}) = 1 {-} 0.25 = 0.75\$.} \\
\texttt{    The event "not *jyka* and not *kwox*" is a subset of "not *jyka*".} \\
\texttt{    \$P(\textbackslash{}neg \textbackslash{}text\{kwox\}) = P(\textbackslash{}neg \textbackslash{}text\{jyka\}) {-} P(\textbackslash{}neg \textbackslash{}text\{jyka\} \textbackslash{}cap \textbackslash{}neg \textbackslash{}text\{kwox\})\$} \\
\texttt{    \$P(\textbackslash{}neg \textbackslash{}text\{kwox\}) = 0.75 {-} 0.71 = 0.04\$.} \\
\\ 
\texttt{    Therefore, the marginal probability of *kwox* is:} \\
\texttt{    \$P(\textbackslash{}text\{kwox\}) = 1 {-} P(\textbackslash{}neg \textbackslash{}text\{kwox\}) = 1 {-} 0.04 = 0.96\$.} \\
\\ 
\texttt{2.  **Calculate the conditional probability of *kwox* given *jyka* (\$P(\textbackslash{}text\{kwox\} | \textbackslash{}text\{jyka\})\$):**} \\
\texttt{    Using the definition of conditional probability:} \\
\texttt{    \$P(\textbackslash{}text\{kwox\} | \textbackslash{}text\{jyka\}) = \textbackslash{}frac\{P(\textbackslash{}text\{jyka\} \textbackslash{}cap \textbackslash{}text\{kwox\})\}\{P(\textbackslash{}text\{jyka\})\}\$} \\
\\ 
\texttt{    Substituting the given values:} \\
\texttt{    \$P(\textbackslash{}text\{kwox\} | \textbackslash{}text\{jyka\}) = \textbackslash{}frac\{0.24\}\{0.25\} = 0.96\$.} \\
\\ 
\texttt{3.  **Comparison:**} \\
\texttt{    \$P(\textbackslash{}text\{kwox\}) = 0.96\$} \\
\texttt{    \$P(\textbackslash{}text\{kwox\} | \textbackslash{}text\{jyka\}) = 0.96\$} \\
\\ 
\texttt{    Since \$P(\textbackslash{}text\{kwox\} | \textbackslash{}text\{jyka\})\$ is equal to \$P(\textbackslash{}text\{kwox\})\$, observing *jyka* does not change the probability of *kwox*. The variables are statistically independent in this specific hypothetical scenario, despite the stated causal direction.} \\
\\ 
\texttt{    The question asks if the chance of *kwox* is **smaller** when observing *jyka*. Since the chance remains exactly the same, the answer is **No**.} \\
            }
        }
    }
    & & \\ \\

    \theutterance \stepcounter{utterance}  
    & & & \multicolumn{2}{p{0.3\linewidth}}{
        \cellcolor[rgb]{0.9,0.9,0.9}{
            \makecell[{{p{\linewidth}}}]{
                \texttt{\tiny{[GM$|$GM]}}
                \texttt{Yes} \\
            }
        }
    }
    & & \\ \\

    \theutterance \stepcounter{utterance}  
    & & & \multicolumn{2}{p{0.3\linewidth}}{
        \cellcolor[rgb]{0.9,0.9,0.9}{
            \makecell[{{p{\linewidth}}}]{
                \texttt{\tiny{[GM$|$GM]}}
                \texttt{no} \\
            }
        }
    }
    & & \\ \\

    \theutterance \stepcounter{utterance}  
    & & & \multicolumn{2}{p{0.3\linewidth}}{
        \cellcolor[rgb]{0.9,0.9,0.9}{
            \makecell[{{p{\linewidth}}}]{
                \texttt{\tiny{[GM$|$GM]}}
                \texttt{game\_result = LOSE} \\
            }
        }
    }
    & & \\ \\

\end{supertabular}
}

\end{document}
